\documentclass[11pt,a4paper]{article}
\usepackage[margin=25mm]{geometry}
\usepackage{amsmath,amssymb}
\usepackage{booktabs}
\usepackage{hyperref}

\title{Single-Element Rigid and Affine Deformation Fields}
\author{}
\date{}

\begin{document}
\maketitle

\section{Definitions}

The cases are generated by \texttt{data/gen\_plate\_data.py}.  The reference
configuration is a square plate centred at the origin, with side length
\begin{equation}
  P = 260.
\end{equation}
The imaging region of interest (ROI) has side length $R=256$ and the generator
uses $N=256$ camera pixels on each axis.  Therefore the physical pixel width is
\begin{equation}
  \Delta = \frac{R}{N} = 1.
\end{equation}
There are eleven frames, indexed by $k\in\{0,\ldots,10\}$.  The common scalar
displacement magnitude at frame $k$ is
\begin{equation}
  s_k = 0.1k\,\Delta = 0.1k.
\end{equation}
All prescribed out-of-plane displacements are zero.  Let the reference
in-plane coordinate be $\mathbf{X}=(X,Y)^\mathsf{T}$ and the deformed coordinate
be $\mathbf{x}=\mathbf{X}+\mathbf{u}(\mathbf{X})$.

The generated nodal coordinates are the corners, edge midpoints, and (where
applicable) centre of this plate.  Since both fields below are at most affine
in $\mathbf{X}$, assigning their exact values at those nodes permits the
standard element interpolation to reproduce the prescribed field exactly.

\section{Rigid translation}

At every node and every point in the plate, the prescribed displacement is
\begin{equation}
  \mathbf{u}_{\mathrm{rigid}}(\mathbf{X},k)
  = \begin{bmatrix}s_k\\s_k\end{bmatrix}.
\end{equation}
Thus the deformation map is a translation,
\begin{equation}
  \mathbf{x}
  = \begin{bmatrix}X+s_k\\Y+s_k\end{bmatrix}.
\end{equation}
The deformation gradient and Jacobian are consequently
\begin{equation}
  \mathbf{F}_{\mathrm{rigid}}=\mathbf{I},
  \qquad J_{\mathrm{rigid}}=\det\mathbf{F}_{\mathrm{rigid}}=1,
\end{equation}
so this case contains neither strain nor rotation.  It is a useful check that
renderers and texture-coordinate mappings move the image without changing its
local shape.

\section{Affine deformation}

The generator assigns the frame-$k$ displacement field
\begin{equation}
  u_X(\mathbf{X},k)=\frac{X+Y}{R}s_k,
  \qquad
  u_Y(\mathbf{X},k)=\frac{Y}{R}s_k.
\end{equation}
Writing $a_k=s_k/R$, the deformation map is
\begin{equation}
  \mathbf{x}
  = \begin{bmatrix}
      (1+a_k)X+a_kY\\
      (1+a_k)Y
    \end{bmatrix}
  = \underbrace{\begin{bmatrix}
      1+a_k & a_k\\
      0 & 1+a_k
    \end{bmatrix}}_{\mathbf{F}_{\mathrm{affine}}}
    \mathbf{X}.
\end{equation}
The displacement gradient is
\begin{equation}
  \nabla\mathbf{u}
  = \begin{bmatrix}a_k & a_k\\0 & a_k\end{bmatrix},
\end{equation}
and the small-strain tensor is
\begin{equation}
  \boldsymbol{\varepsilon}
  = \frac{1}{2}\left(\nabla\mathbf{u}+\nabla\mathbf{u}^{\mathsf T}\right)
  = \begin{bmatrix}a_k & a_k/2\\a_k/2 & a_k\end{bmatrix}.
\end{equation}
Hence the field combines equal normal strains in $X$ and $Y$ with engineering
shear strain $\gamma_{XY}=a_k$.  Its area change is
\begin{equation}
  J_{\mathrm{affine}}=\det\mathbf{F}_{\mathrm{affine}}=(1+a_k)^2.
\end{equation}
For the largest generated displacement, $s_{10}=1$, so
$a_{10}=1/256\approx3.90625\times10^{-3}$: this is a deliberately small,
spatially uniform affine deformation.

\section{Reference UV coordinates}

For completeness, reference UVs are stored with the coordinate-to-ROI mapping
\begin{equation}
  U=\frac{X+R/2}{R},
  \qquad
  V=\frac{Y+R/2}{R}.
\end{equation}
Thus the ROI $[-R/2,R/2]^2$ maps to $[0,1]^2$.  The plate is slightly larger
than the ROI ($P>R$), so corner-node UVs lie marginally outside the unit square;
the renderer texture setup supplies the required padded texture coverage.

\begin{table}[h]
\centering
\begin{tabular}{@{}lll@{}}
\toprule
Case & Displacement $\mathbf{u}(\mathbf{X},k)$ & Map $\mathbf{x}(\mathbf{X},k)$ \\
\midrule
Rigid & $(s_k,s_k)^\mathsf{T}$ & $(X+s_k,Y+s_k)^\mathsf{T}$ \\
Affine & $a_k(X+Y,Y)^\mathsf{T}$ & $((1+a_k)X+a_kY,(1+a_k)Y)^\mathsf{T}$ \\
\bottomrule
\end{tabular}
\caption{Deformation fields written in terms of $s_k=0.1k$ and $a_k=s_k/R$.}
\end{table}

\end{document}
